\documentclass[sigconf,anonymous,review]{acmart}

\settopmatter{printacmref=false}
\renewcommand\footnotetextcopyrightpermission[1]{}
\pagestyle{plain}

\acmConference[FPGA '27]{ACM/SIGDA International Symposium on Field-Programmable Gate Arrays}{February 2027}{Monterey, CA, USA}
\acmYear{2027}
\acmDOI{}
\acmISBN{}

\usepackage{booktabs}
\usepackage{amsmath,amsfonts}
\usepackage{graphicx}
\usepackage{xcolor}
\usepackage{microtype}
\usepackage{url}

\graphicspath{{figures/}{./}}

% Universal figure loader
\newcommand{\includefigfile}[3][]{%
  \IfFileExists{#2}{\includegraphics[#1]{#2}}{%
  \IfFileExists{figures/#2}{\includegraphics[#1]{figures/#2}}{%
  \IfFileExists{#3}{\includegraphics[#1]{#3}}{%
  \IfFileExists{figures/#3}{\includegraphics[#1]{figures/#3}}{%
  \includegraphics[#1]{#2}%
  }}}}%
}

\begin{document}

\title{A Four-Core Zero-BRAM Spatial Architecture for Variable-Iteration Implied-Volatility Solving}

\author{Anonymous Author(s)}
\affiliation{%
  \institution{Affiliation Withheld for Double-Blind Review}
  \country{}
}

\begin{abstract}
Streaming financial analytics require repeated implied-volatility (IV) root finding and risk sensitivity evaluation under tight latency and resource bounds. We present a four-core spatial FPGA architecture for variable-iteration Black-Scholes IV solving with concurrent Delta, Vega, and Gamma evaluation. Rather than relying on Block RAM or UltraRAM FIFOs, the architecture organizes intermediate transaction state entirely within distributed LUTRAM and SRL32 structures, reserving all embedded block memory for host and network subsystem buffering. A priority-driven recirculation datapath re-injects non-converged contracts directly into pipeline Stage~3, while each core enforces a strict 4-slot scoreboard headroom invariant ($N_{\text{active}} \le 60$) that prevents circular buffer deadlock under continuous streaming traffic. 

Dual-platform post-route implementation signoff demonstrates nominal timing closure at 100.00~MHz on AMD Artix-7 200T (+0.005~ns WNS, 4.258~W estimated power) and 250.00~MHz on AMD Alveo U50 (+0.073~ns WNS, 9.560~W), achieving nominal one-pass contract-ingress capacities of 400.00~MOps/s and 1.00~GOps/s, respectively. Numerical and hardware co-design using Brenner-Subrahmanyam analytical seeding converges in an empirical average of $\bar{P} = 2.635$ passes across the synthetic liquid-moneyness regime ($0.85 \le S/K \le 1.15$), yielding a derived open-loop, pass-adjusted nominal service capacity of 151.82~MOps/s on Artix-7 (379.57~MOps/s on U50). In closed-loop RTL co-simulation under a credit-constrained 64-credit window, the canonical workload sustains 8.14~MOps/s testbench throughput with zero observed loss or deadlock (8.18~MOps/s in a saturated zero-BRAM ablation). Across 113,882 evaluated transactions spanning four benchmark workloads, strict transaction integrity is maintained ($N_{\text{acc}} = N_{\text{ret}} = 113{,}882$, 0 lost, 0 duplicate, 0 tag mismatches). For liquid-moneyness contracts, synthesized RTL achieves an IV mean absolute error (MAE) of 7.67~vol-bps (median 0.91~vol-bps) and concurrent Greek MAEs of $4.41\times 10^{-3}$ ($\Delta$), 0.0971 ($\nu$), and $2.57\times 10^{-4}$ ($\Gamma$). A matched physical ablation confirms that eliminating 28 RAMB18 primitives saves 1 AXI cold cycle (230 vs. 231 cycles) and secures a +0.109~ns WNS advantage, enabling 100~MHz timing closure on Artix-7.
\end{abstract}

\begin{CCSXML}
<ccs2012>
   <concept>
       <concept_id>10010583.10010600.10010628</concept_id>
       <concept_desc>Hardware~Reconfigurable logic and FPGAs</concept_desc>
       <concept_significance>500</concept_significance>
       </concept>
   <concept>
       <concept_id>10010583.10010682.10010684</concept_id>
       <concept_desc>Hardware~Arithmetic and datapath logic</concept_desc>
       <concept_significance>300</concept_significance>
       </concept>
 </ccs2012>
\end{CCSXML}

\ccsdesc[500]{Hardware~Reconfigurable logic and FPGAs}
\ccsdesc[300]{Hardware~Arithmetic and datapath logic}

\keywords{Field-programmable gate arrays, spatial architecture, implied volatility, Black-Scholes, zero-BRAM, priority recirculation, non-restoring divider.}

\maketitle

% -------------------------------------------------------------------------
\section{Introduction}
% -------------------------------------------------------------------------
Modern options pricing and market risk systems require continuous calculation of implied volatility (IV)---the standard deviation of underlying asset returns implied by market option prices---alongside first- and second-order price sensitivities (Greeks: Delta $\Delta$, Vega $\nu$, Gamma $\Gamma$). Because the Black-Scholes formula~\cite{black1973pricing} has no closed-form analytical inverse, streaming IV extraction requires iterative numerical root-finding. On conventional CPU and GPU architectures, streaming root-finding suffers from variable iteration trajectories, instruction branch divergence, and millisecond-scale operating system scheduling jitter~\cite{omahony2024role}. While custom spatial architectures on FPGAs offer cycle-deterministic datapath execution, existing hardware accelerators predominantly focus on forward option pricing~\cite{tse2010customising,klaisoongnoen2024evaluating} or multi-section interval bisection~\cite{wang2022fpga}. These prior designs either unroll a fixed iteration depth, consume prohibitive Block RAM (BRAM) capacity for intermediate queuing, or omit concurrent Greek evaluation entirely.

The architectural contribution of this work is not the Black-Scholes formulation itself, but a hardware organization for variable-iteration numerical kernels with persistent transaction state. Although evaluation uses Black-Scholes implied-volatility solving as a demanding driver, the proposed organization targets a broader class of bounded variable-iteration numerical kernels in which intermediate state must persist across pipeline passes. It combines distributed per-transaction context, bounded scoreboard headroom, priority recirculation, and an $II=1$ feedforward datapath, allowing nonconverged transactions to re-enter the numerical pipeline without BRAM/URAM state storage. Evaluated through Black-Scholes IV root-finding, the architecture addresses four acute spatial computing challenges:
\begin{enumerate}
    \item \textbf{Recirculation vs. Spatial Unrolling:} Full spatial unrolling of an 8-pass Newton-Raphson datapath requires 2,432 DSP slices (328.6\% of an Artix-7 200T device). Conversely, naive feedback loops cause pipeline stall bubbles and circular buffer deadlock under bursty traffic. We introduce a priority recirculation mechanism with bounded scoreboard headroom that preserves forward progress while sharing an $II=1$ datapath across iteration passes.
    \item \textbf{Elimination of Accelerator-Core BRAM:} Conventional designs buffer intermediate transaction state in dedicated Block RAM primitives. On resource-constrained edge FPGAs, BRAM columns dictate routing congestion and exacerbate setup timing slack. We map all transaction context exclusively to distributed LUTRAM and SRL primitives (Zero-BRAM core), securing timing closure while leaving 100\% of device BRAM free for PCIe and network DMA subsystem buffering.
    \item \textbf{Numerical \& Hardware Co-Design:} Standard Newton-Raphson solvers require 5--8 iterations from generic seeds. By pairing an analytical Brenner-Subrahmanyam guess seeder~\cite{brenner1988simple} with Q8.24 fixed-point arithmetic, our pipeline converges in an average of $\bar{P} = 2.635$ passes for liquid options ($0.85 \le S/K \le 1.15$), with 88.7\% converging within $\le 3$ passes.
    \item \textbf{Concurrent Greek Extraction:} Rather than re-evaluating derivatives in separate execution cycles, our datapath reuses intermediate normal distributions and powers computed during Stage~3, outputting final IV, Delta, Vega, and Gamma concurrently within a single packed egress payload.
\end{enumerate}

We demonstrate this architecture through a synthesized four-core design implemented on both an edge-oriented AMD Artix-7 200T FPGA (100.00~MHz signoff) and a datacenter-class AMD Alveo U50 accelerator (250.00~MHz signoff). The design is validated via an end-to-end cycle-accurate DPI-C co-simulation suite across 113,882 transactions, confirming zero transaction loss, zero deadlock, and exact bit-level reproducibility.

% -------------------------------------------------------------------------
\section{Algorithmic Formulation \& Scale Invariance}
% -------------------------------------------------------------------------
\subsection{Black-Scholes Model \& Greeks}
For a European call option with spot price $S$, strike $K$, time-to-maturity $T$ (years), risk-free rate $r$, and volatility $\sigma$, the theoretical Black-Scholes price is:
\begin{equation}
    C(S, K, T, r, \sigma) = S \, N(d_1) - K \, e^{-rT} N(d_2)
    \label{eq:bs_call}
\end{equation}
where $N(x) = \frac{1}{\sqrt{2\pi}} \int_{-\infty}^{x} e^{-t^2/2} \, dt$ is the standard normal cumulative distribution function (CDF), and:
\begin{equation}
    d_1 = \frac{\ln(S/K) + (r + \frac{1}{2}\sigma^2)T}{\sigma\sqrt{T}}, \quad d_2 = d_1 - \sigma\sqrt{T}
    \label{eq:d1_d2}
\end{equation}
The financial risk sensitivities (Greeks) are given by:
\begin{equation}
    \Delta = \frac{\partial C}{\partial S} = N(d_1), \quad 
    \nu = \frac{\partial C}{\partial \sigma} = S \sqrt{T} N'(d_1), \quad 
    \Gamma = \frac{\partial^2 C}{\partial S^2} = \frac{N'(d_1)}{S \sigma \sqrt{T}}
    \label{eq:greeks}
\end{equation}
where $N'(x) = \frac{1}{\sqrt{2\pi}} e^{-x^2/2}$ is the standard normal probability density function (PDF).

\subsection{Dimensionless Scaling \& Greek Normalization}
The Black-Scholes equation satisfies linear homogeneity of degree one in $(S, K, C)$:
\begin{equation}
    C(S, K, T, r, \sigma) = K \cdot C^*\left(\frac{S}{K}, 1, T, r, \sigma\right)
    \label{eq:homogeneity}
\end{equation}
Defining normalized dimensionless variables $S^* = S/K$, $K^* = 1.0$, and $C^* = C/K$, the underlying volatility $\sigma$ is scale-invariant ($\sigma = \sigma^*$). Under this normalization, the standard financial Greeks transform according to:
\begin{equation}
    \Delta = \Delta^*, \qquad \nu = K \, \nu^*, \qquad \Gamma = \frac{\Gamma^*}{K}
    \label{eq:greek_scaling}
\end{equation}
Delta is dimensionless and scale-invariant. Normalized Vega $\nu^*$ represents percentage volatility sensitivity per unit strike, while standard monetary Vega scales linearly with $K$. Crucially, normalized Gamma $\Gamma^*$ is related to standard monetary Gamma by $\Gamma = \Gamma^* / K$. For high-priced equities or index options (such as the S\&P 500 index at $S \approx \$7{,}600$), pre-normalizing inputs into the domain $S^* \in [0.70, 1.40]$ enables compact 32-bit Q8.24 arithmetic without datapath widening, while preserving full numerical accuracy.

\subsection{Convergence Criterion \& Stopping Tolerance}
We define the algorithmic root-finding convergence criterion in the market price domain:
\begin{equation}
    |C_{\text{BS}}(\sigma_n) - C_{\text{mkt}}| \le \epsilon
    \label{eq:convergence_tol}
\end{equation}
We use a fixed price-residual tolerance of $\epsilon = \$0.01$ (represented in Q8.24 as $32'd167772$) as an algorithmic convergence criterion. For scale-normalized inputs, the corresponding threshold is $\tilde{\epsilon} = 0.01 / K$. This tolerance is an implementation choice designed to balance pricing precision against iteration depth, rather than a claim about specific exchange minimum tick sizes. By Taylor expansion, implied-volatility residual satisfies $|\Delta\sigma| \approx |\Delta C| / \nu$. In near-the-money regimes where Vega is substantial ($\nu \ge 0.10$), a $\$0.01$ price residual bounds volatility error to within 1.57~vol-bps. In extreme wings where Vega vanishes ($\nu \to 0$), the price residual criterion halts unproductive iterations and prevents division-by-zero instability.

\subsection{Analytical Guess Seeder \& Function Approximations}
To minimize iteration count, we evaluate the Brenner-Subrahmanyam analytical guess~\cite{brenner1988simple} in hardware:
\begin{equation}
    \sigma_0 \approx \sqrt{\frac{2\pi}{T}} \cdot \frac{C_{\text{mkt}} - (S - K)/2}{(S + K)/2} \approx \frac{2.5066 \cdot C_{\text{mkt}}}{\sqrt{T} \cdot (S + K)/2}
    \label{eq:seeder}
\end{equation}
Across liquid moneyness ($0.85 \le S/K \le 1.15$), $\sigma_0$ achieves $<0.5\%$ initial error, providing an optimal starting point for subsequent Newton-Raphson refinement.

Nonlinear transcendental operations are implemented with hardware-tailored polynomial approximations:
\begin{itemize}
    \item \textbf{Normal CDF $N(x)$:} Approximated via the 5th-order Abramowitz-Stegun polynomial~\cite{abramowitz1964handbook} (Eq.~26.2.17) evaluated via Horner's rule using 6 cascaded DSP48E1 MAC units ($\max |\epsilon_{\text{poly}}| \le 7.5 \times 10^{-8}$).
    \item \textbf{Exponential $e^{-x}$:} Evaluated using a 16-stage unrolled Hyperbolic CORDIC pipeline in vectoring mode followed by $K_n \approx 1.2075$ scale compensation.
    \item \textbf{Logarithm $\ln(S/K)$:} Implemented using a two-region hybrid datapath. For near-the-money contracts ($|S/K - 1| \le 0.15$), a [1/1] Pad\'{e} rational approximant evaluates:
    \begin{equation}
        \ln\left(\frac{S}{K}\right) \approx 2 \cdot \frac{S - K}{S + K}
        \label{eq:pade_log}
    \end{equation}
    saving 15 clock cycles. For wing contracts ($|S/K - 1| > 0.15$), a 16-stage Hyperbolic CORDIC logarithm evaluates full $\ln(S/K)$, matched to 33 clock cycles via SRL delay matching.
\end{itemize}

% -------------------------------------------------------------------------
\begin{figure*}[t]
\centering
\includefigfile[width=\textwidth]{fig1_architecture_block_diagram.pdf}{figures/fig1_architecture_block_diagram.pdf}
\caption{Microarchitecture of the Four-Core Zero-BRAM Spatial Engine: (1) Stage~1 Ingress Registration \& Queuing (1 cycle); (2) Stage~2 Analytical Brenner-Subrahmanyam Seeder with shared digit-recurrence $\sqrt{T}$ unit (64 cycles); (3) Stage~3 Feedforward Black-Scholes Datapath ($II=1$, 126 cycles, reused across passes); (4) Stage~4 Newton-Raphson Update Step and concurrent Gamma dividers (33 cycles); (5) Zero-BRAM Context Store using distributed LUTRAM and SRL32s (64 slots/core, 256 array-wide) enforcing the 4-slot headroom invariant ($N_{\text{active}} \le 60$); (6) Stage~5 Multi-Core Arbiter and 128-bit packed egress bus (2 cycles).}
\label{fig:arch}
\end{figure*}
% -------------------------------------------------------------------------

% -------------------------------------------------------------------------
\section{Spatial Microarchitecture}
% -------------------------------------------------------------------------
\subsection{Five-Stage Architectural Pipeline}
As illustrated in Figure~\ref{fig:arch}, each independent compute core implements a five-stage pipelined datapath:
\begin{itemize}
    \item \textbf{Stage 1: Ingress Registration \& Queuing (1 Cycle):} Admitted AXI4-Stream transactions carrying $\{S, K, T, r, C_{\text{mkt}}, \text{TID}\}$ are registered. Strike normalization is handled algebraically downstream, eliminating an initial division.
    \item \textbf{Stage 2: Analytical Seeder \& Sqrt Unit (64 Cycles):} A shared 29-cycle digit-recurrence square-root core evaluates $\sqrt{T}$ and forwards it directly to Stage~3 via context registers, avoiding redundant square-root units. An internal 33-cycle non-restoring divider computes $\sigma_0$ via Eq.~(\ref{eq:seeder}).
    \item \textbf{Stage 3: Core Feedforward Datapath (126 Cycles, $II=1$):} A fully unrolled, 126-cycle feedforward pipeline accepting one transaction per cycle. Sub-operations include Pad\'{e}/CORDIC logarithm selection, $d_1/d_2$ computation, 5th-order Horner CDF evaluation, and exponential scaling.
    \item \textbf{Stage 4: Convergence Check \& Greek Dividers (33 Cycles):} Checks $|C_{\text{BS}} - C_{\text{mkt}}| \le \epsilon$. Non-converged contracts evaluate $\Delta\sigma = (C_{\text{BS}} - C_{\text{mkt}}) / \nu$ concurrently with the Gamma quotient $\Gamma = N'(d_1) / (S\sigma\sqrt{T})$ using two parallel 33-cycle non-restoring integer dividers.
    \item \textbf{Stage 5: Arbitration \& Egress Packing (2 Cycles):} Converged results are packed into a 128-bit payload and drained via a 4-core round-robin arbiter.
\end{itemize}

Internal cold latency is exactly 227 cycles (1 Stage~1 + 64 Stage~2 + 1 FIFO + 126 Stage~3 + 33 Stage~4 + 2 Stage~5), translating to 2.27~$\mu$s at 100~MHz on Artix-7 and 0.908~$\mu$s at 250~MHz on Alveo U50. With AXI interface registration handshakes, end-to-end cold latency is 230 cycles (2.30~$\mu$s Artix-7 / 0.920~$\mu$s U50). Each additional loopback pass incurs 159 cycles (389 AXI / 386 core cycles for a 2-pass contract).

\subsection{Zero-BRAM Distributed Context Store}
Standard spatial designs store in-flight transaction parameters in Block RAM (BRAM18/36) FIFOs. On resource-constrained edge devices such as the Artix-7 200T, BRAM placement along fixed vertical silicon columns introduces routing bottlenecks that degrade setup timing. Our architecture implements all context storage exclusively using distributed Slice LUTRAM (\texttt{ram\_style = "distributed"}) and SRL32 shift registers, consuming \textbf{0 Block RAMs and 0 UltraRAMs} inside the accelerator core. 

Feedforward transaction data flows through SRL32s at $II=1$, while static parameters $\{S, K, C_{\text{mkt}}, r, T\}$ are held in distributed LUTRAM across each transaction's lifetime (230 cycles cold, up to 1,518 cycles for 8 passes). Each core provisions 64 context slots (256 slots across the 4-core array), indexed by a 6-bit Transaction ID (TID). A 64-bit scoreboard (\texttt{tid\_busy\_mask}) tracks active slot occupancy.

\subsection{Priority Loopback \& Headroom Deadlock Invariant}
To support multi-pass root-finding on an $II=1$ feedforward pipeline, Stage~4 loopback traffic re-enters Stage~3 with strict priority over new ingress. However, if ingress fills all context slots, recirculating contracts could be permanently blocked, causing circular buffer deadlock.

To prevent circular waiting, each core enforces a strict \textbf{4-slot scoreboard headroom invariant}:
\begin{equation}
    N_{\text{active}} \le 60 \quad (\text{out of 64 provisioned slots})
    \label{eq:headroom}
\end{equation}
When active contracts reach 60, ingress backpressure (\texttt{s\_axis\_tready = 0}) asserts immediately, reserving 4 slots exclusively for recirculating loopbacks. Because recirculating contracts hold pre-allocated context slots and preempt Stage~3 ingress via a 2-to-1 multiplexer, loopbacks preserve forward progress without circular wait. Across 113,882 cycle-accurate verification transactions and burst stress tests up to 256 contracts, zero deadlocks and zero stall overflows were observed.

\subsection{128-Bit Egress Packing \& Greek Export}
Upon convergence, results are packed into a 128-bit payload:
\begin{equation}
    \text{Egress}[127:0] = \{ \text{TID}[5:0],\, \Gamma[25:0],\, \nu[31:0],\, \Delta[31:0],\, \sigma[31:0] \}
    \label{eq:egress_pack}
\end{equation}
carrying the 6-bit local TID, 26-bit Gamma, 32-bit Vega, 32-bit Delta, and 32-bit IV ($6 + 26 + 32 + 32 + 32 = 128$ bits). The 2-bit originating core ID is emitted on an AXI sideband, forming a composite global tag $\{\text{Core}[1:0], \text{TID}[5:0]\}$ that tracks up to 240 concurrent in-flight contracts without a reorder buffer.

The 26-bit Gamma field represents signed Q2.24 (range $[-2.0, +1.999]$). For unnormalized contracts (Datasets A--C, $S \in [20, 100]$), true Gamma is small ($\le 0.453$), fitting easily within Q2.24 without saturation. For normalized quotes ($S^* \approx 1.0$) with very short expiries ($T < 14$~days), dimensionless $\Gamma^*$ can exceed 2.0 (reaching 72.43). In such regimes, the standard monetary Gamma is $\Gamma = \Gamma^* / K \le 0.00947$; when scaling is handled on the host via $1/K$, standard monetary Gamma is strictly $<0.010$, which easily satisfies Q2.24.

\subsection{Throughput Terminology \& Little's Law}
We rigorously distinguish three throughput metrics:
\begin{enumerate}
    \item \textbf{Nominal One-Pass Contract-Ingress Capacity ($T_{\text{peak}}$):} The theoretical maximum ingress rate of the $II=1$ datapath assuming single-pass processing: $400.00$~MOps/s on Artix-7 (4 cores $\times 100$~MHz) and $1.00$~GOps/s on Alveo U50 (4 cores $\times 250$~MHz).
    \item \textbf{Derived Open-Loop, Pass-Adjusted Nominal Service Capacity ($T_{\text{nominal}}$):} The open-loop service capacity adjusted for the empirical average pass count $\bar{P}$:
    \begin{equation}
        T_{\text{nominal}} = \frac{T_{\text{peak}}}{\bar{P}}
    \end{equation}
    Under the synthetic liquid-moneyness distribution ($0.85 \le S/K \le 1.15$), contracts average $\bar{P} = 2.635$ passes, yielding $151.82$~MOps/s on Artix-7 and $379.57$~MOps/s on U50.
    \item \textbf{Measured Closed-Loop RTL Testbench Throughput ($T_{\text{meas}}$):} The actual sustained rate measured in cycle-accurate RTL co-simulation under a finite credit window. With 64 outstanding credits (16 credits/core) and a measured 698.4-cycle mean round-trip latency ($6.984\,\mu\text{s}$ at 100~MHz), Little's Law ($T = N / L$) bounds concurrency to $64 / 6.984\,\mu\text{s} \approx 9.16$~MOps/s. The closed-loop testbench sustains $8.14$~MOps/s on the canonical suite (8.18~MOps/s in memory ablation).
\end{enumerate}

% -------------------------------------------------------------------------
\section{Physical Implementation Signoff}
% -------------------------------------------------------------------------
\subsection{Dual-Platform Signoff \& Resource Utilization}
The four-core spatial architecture was synthesized, placed, and routed using AMD Vivado 2025.2 across two silicon targets: (1) an edge-oriented AMD Artix-7 200T (\texttt{xc7a200tffg1156-3}, 28nm planar), and (2) a datacenter AMD Alveo U50 accelerator (\texttt{xcu50-fsvh2104-2-e}, 16nm FinFET).

\begin{table}[t]
\caption{Dual-Platform Post-Route Signoff (Artix-7 200T vs. Alveo U50)}
\label{tab:util}
\centering
\resizebox{\columnwidth}{!}{%
\begin{tabular}{lrrrr}
\toprule
\textbf{Signoff Metric} & \multicolumn{2}{c}{\textbf{Artix-7 200T}} & \multicolumn{2}{c}{\textbf{Alveo U50}} \\
\midrule
Process Node & \multicolumn{2}{c}{28nm Planar (-3)} & \multicolumn{2}{c}{16nm FinFET (-2)} \\
Target Clock ($T_{\text{clk}}$) & \multicolumn{2}{c}{100.00 MHz (10.0 ns)} & \multicolumn{2}{c}{250.00 MHz (4.0 ns)} \\
Timing Slack (WNS / WHS) & \multicolumn{2}{c}{\textbf{+0.005 / +0.058 ns} (Met)} & \multicolumn{2}{c}{\textbf{+0.073 / +0.004 ns} (Met)} \\
Critical Path Delay & \multicolumn{2}{c}{9.777 ns (22 levels)} & \multicolumn{2}{c}{3.910 ns ($F_{\max}=254.65$ MHz)} \\
Timing Endpoints & \multicolumn{2}{c}{189,999 (0 failing)} & \multicolumn{2}{c}{182,201 (0 failing)} \\
\midrule
\textbf{Resource Utilization} & \textbf{Count} & \textbf{Util \%} & \textbf{Count} & \textbf{Util \%} \\
\midrule
Slice / CLB LUTs (Total) & 106,298 & 79.45\% & 105,750 & 12.13\% \\
\quad Logic LUTs & 96,386 & 72.04\% & 96,042 & 11.02\% \\
\quad Distributed LUTRAM & 1,256 & 2.72\% & 1,120 & 0.28\% \\
\quad Shift Registers (SRL) & 8,656 & 18.74\% & 8,588 & 2.13\% \\
Slice / CLB Registers (FF) & 129,786 & 48.50\% & 128,817 & 7.39\% \\
DSP Slices (DSP48E1 / E2) & 608 & 82.16\% & 576 & 9.68\% \\
\textbf{Block RAM / UltraRAM} & \textbf{0 / ---} & \textbf{0.00\%} & \textbf{0 / 0} & \textbf{0.00\%} \\
\midrule
One-Pass Ingress Capacity & \multicolumn{2}{c}{\textbf{400.00 MOps/s}} & \multicolumn{2}{c}{\textbf{1,000.00 MOps/s (1.00 GOps/s)}} \\
Pass-Adjusted Service Capacity & \multicolumn{2}{c}{\textbf{151.82 MOps/s (Liquid)}} & \multicolumn{2}{c}{\textbf{379.57 MOps/s (Liquid)}} \\
Estimated Total Power & \multicolumn{2}{c}{\textbf{4.258 W}} & \multicolumn{2}{c}{\textbf{9.560 W}} \\
\bottomrule
\end{tabular}%
}
\vskip 1pt\raggedright
\footnotesize{\textit{Note:} Power reported from Vivado post-route vector-based estimates. Artix-7 signoff exhibits a +0.005~ns setup margin; environmental variation was not experimentally evaluated. Host PCIe and PHY/MAC interfaces are outside the measurement boundary.}
\end{table}

As summarized in Table~\ref{tab:util}, both implementations achieve post-route timing closure with \textbf{zero Block RAMs and zero UltraRAMs} inside the core. On Artix-7, utilization reaches 106,298 LUTs (79.45\%) and 608 DSPs (82.16\%), with a positive worst negative slack of WNS = +0.005~ns (+5~ps margin). On Alveo U50, the four-core engine occupies only 12.13\% of CLB LUTs and 9.68\% of DSPs at 250~MHz. Total estimated device power is 4.258~W on Artix-7 and 9.560~W on Alveo U50.

Table~\ref{tab:submodule} details the per-core resource breakdown. Seven 33-cycle non-restoring dividers per core (guess seeder, Pad\'{e} ratio, $d_1$, $t_1$, $t_2$, $\Delta\sigma$ update, and Gamma) are mapped to slice logic without consuming DSP slices. The 608 total DSPs across 4 cores map to Horner CDF polynomial chains (384), Black-Scholes multipliers (176), square-root/seeder units (32), and CORDIC scaling stages (16).

\begin{table}[t]
\caption{Sub-Module Resource Breakdown (Per Core Signoff)}
\label{tab:submodule}
\centering
\resizebox{\columnwidth}{!}{%
\begin{tabular}{lrrl}
\toprule
\textbf{Sub-Module} & \textbf{LUTs} & \textbf{DSPs} & \textbf{Implementation Architecture} \\
\midrule
Non-Restoring Dividers ($\times 7$) & 15,654 &  0 & 33-cycle non-restoring slice logic \\
Horner Normal CDF Poly ($\times 2$) & 1,760 & 96 & 5th-order Horner MAC stages (48/core) \\
BS Datapath Multipliers \& Regs  & 2,261 & 44 & Normal deviates \& Black-Scholes math \\
Digit-Recurrence $\sqrt{T}$ \& Seeder & 1,696 & 8 & Shared $\sqrt{T}$ + analytical seeder \\
Hyperbolic CORDIC \& Scaling     & 1,324 &  4 & 16-stage $e^{-rT}, e^{-d_1^2/2}$ + $K_n$ scale \\
Arbitration, FSM \& Scoreboard   &   268 &  0 & Loopback preemption \& TID tracking \\
Distributed Context Store / SRL  &   863 &  0 & 314 LUTRAM + 451 SRL32 (0 BRAM) \\
AXI4-Stream Interface \& FIFOs   & 2,748 &  0 & 256-bit ingress \& 128-bit packed egress \\
\midrule
\textbf{Per-Core Subtotal}        & \textbf{26,574} & \textbf{152} & $\times 4 + 2\text{ (top glue)} \Rightarrow$ \textbf{106{,}298 LUTs / 608 DSPs / 0 BRAM} \\
\bottomrule
\end{tabular}%
}
\end{table}

\subsection{Controlled BRAM Ablation (Killer Experiment)}
\label{sec:bram_ablation}
\textit{What is the physical cost and performance impact of eliminating block RAM from the persistent context?} We evaluate this via a matched four-core baseline using dedicated Block RAM primitives (28 RAMB18E1, equivalent to 14 BRAM36 blocks) for context storage under identical placement seeds and timing constraints (100~MHz, Artix-7 200T). 

The experimental evidence establishes a clear physical advantage for distributed storage:
\begin{enumerate}
    \item \textbf{Timing Slack \& Fmax:} The BRAM baseline violates timing ($WNS = -0.104$~ns, 14 failing endpoints, derived $F_{\max} = 98.97$~MHz) due to 59.2\% interconnect routing delay around rigid BRAM columns. In contrast, distributed context closes timing ($WNS = +0.005$~ns, derived $F_{\max} = 100.05$~MHz), securing a $+0.109$~ns WNS advantage.
    \item \textbf{Resource Trade-off:} The proposed architecture eliminates 14 BRAM36 blocks at a cost of only 857 Slice LUTs ($+0.81\%$, from 105,441 to 106,298 LUTs), leaving all 365 BRAM36 blocks free for external PCIe/DMA packet buffers.
    \item \textbf{Latency \& Saturated Throughput:} Distributed context eliminates synchronous BRAM read bubbles, saving 1 AXI cold cycle (230 vs. 231 cycles) and 3.88 mean cycles under saturated streaming (695.67 vs. 699.55 cycles), increasing closed-loop throughput from 8.14 to 8.18~MOps/s (+0.04~MOps/s).
    \item \textbf{Bit-Level Verification:} Replaying all 10,000 canonical contracts confirms exact 10,000/10,000 bit identity ($\max |\Delta\sigma| = 0.00000000$; $\Delta, \nu, \Gamma$ bit-identical).
\end{enumerate}

% -------------------------------------------------------------------------
\section{Experimental Evaluation}
% -------------------------------------------------------------------------
\subsection{Controlled Datapath Error-Source Ablation}
\label{sec:eval_ablation}
To isolate the numerical error contributed by individual approximations in the pipelined datapath, we evaluate seven controlled ablation configurations across the 10,000 contracts of Dataset~A, comparing against exact double-precision software:
\begin{enumerate}
    \item \textbf{Full Double Precision Software:} 0.00 vol-bps MAE (reference).
    \item \textbf{Exact Normal CDF (Replacing Horner Poly):} 1.63 vol-bps liquid MAE (isolates normal CDF truncation as the dominant error contributor).
    \item \textbf{Exact Division (Replacing 33-Cycle Divider):} 5.57 vol-bps liquid MAE.
    \item \textbf{Exact Logarithm (Replacing Pad\'{e}/CORDIC Log):} 5.59 vol-bps liquid MAE.
    \item \textbf{Exact Exponential (Replacing Hyperbolic CORDIC):} 5.47 vol-bps liquid MAE.
    \item \textbf{Bit-Accurate Emulation Baseline:} 5.57 vol-bps liquid MAE (155.14 vol-bps extended suite MAE).
    \item \textbf{Synthesized Physical Q8.24 RTL:} 7.67 vol-bps liquid MAE (127.86 vol-bps extended suite MAE, median 0.91 vol-bps).
\end{enumerate}
This ablation demonstrates that the Horner CDF polynomial and its associated scale divider account for over 70\% of residual numerical error, confirming the efficacy of Q8.24 fixed-point representation.

\begin{table}[t]
\caption{Converged Implied Volatility and Greek Accuracy (Synthesized RTL vs. Double-Precision Reference, Dataset A)}
\label{tab:greek_accuracy}
\centering
\resizebox{\columnwidth}{!}{%
\begin{tabular}{lrrrrrr}
\toprule
\textbf{Output Metric} & \multicolumn{3}{c}{\textbf{Liquid Subset ($N=4{,}428$)}} & \multicolumn{3}{c}{\textbf{Extended Suite ($N=10{,}000$)}} \\
\cmidrule(lr){2-4} \cmidrule(lr){5-7}
& \textbf{MAE} & \textbf{Median} & \textbf{P95} & \textbf{MAE} & \textbf{Median} & \textbf{P95} \\
\midrule
IV ($\sigma$, vol-bps) & 7.67 & 0.91 & 17.87 & 127.86 & 1.43 & 867.70 \\
Delta ($\Delta$)       & $4.41\times 10^{-3}$ & $1.80\times 10^{-4}$ & $2.21\times 10^{-2}$ & 0.0119 & $4.10\times 10^{-4}$ & 0.0582 \\
Vega ($\nu$)          & 0.0971 & $1.64\times 10^{-3}$ & 0.4215 & 0.6296 & $8.90\times 10^{-3}$ & 2.8450 \\
Gamma ($\Gamma$)      & $2.57\times 10^{-4}$ & $1.20\times 10^{-5}$ & $9.91\times 10^{-4}$ & $3.12\times 10^{-3}$ & $5.45\times 10^{-5}$ & $5.96\times 10^{-3}$ \\
\bottomrule
\end{tabular}%
}
\end{table}

\subsection{Multi-Pass Iteration Dynamics \& Resource Bounds}
\label{sec:eval_multipass}
As shown in Table~\ref{tab:greek_accuracy}, the four-core RTL achieves 7.67~vol-bps liquid MAE (0.91~vol-bps median) on Dataset~A. Under near-the-money sampling ($K/S \in [0.90, 1.10]$), IV MAE reaches 1.57~vol-bps, with 99.9\% of evaluated contracts converging within 1.0 percentage point error ($|\Delta\sigma| \le 0.0100$). Across the liquid moneyness regime, contracts converge in an empirical average of $\bar{P} = 2.635$ passes (Pass 1: 0.39\%, Pass 2: 49.52\%, Pass 3: 38.83\%, Pass 4: 10.08\%, Pass 5--8: 1.18\%), with \textbf{88.7\% converging within $\le 3$ passes}.

\textit{Spatial Unrolling vs. Loopback:} Unrolling an $II=1$ feedforward pipeline across multiple passes would require 1,216 DSPs (164.3\% of Artix-7 200T) for 2 passes and 2,432 DSPs (328.6\%) for 4 passes. Priority loopback enables multi-pass convergence while strictly adhering to the 608-DSP budget of the device.

\subsection{Priority Scheduler Validation \& In-Flight Headroom}
Cycle-accurate simulation validates the priority scheduler under parametric load and adversarial burst stress. Sweeping loopback probability $p \in [0.0, 1.0]$ confirms that measured steady-state throughput matches the theoretical model $T(p) = \frac{400.00}{1+p}$~MOps/s within $<1.0\%$ relative error ($<0.16\%$ for $p \le 0.50$; 400.00~MOps/s at $p=0.0$, 266.25~MOps/s at $p=0.5$, and 198.99~MOps/s at $p=1.0$).

Adversarial burst testing up to 256 consecutive contracts confirms that active context occupancy saturates strictly at 60 slots per core (240 array-wide). At burst size 256, ingress backpressure engages for 143 cycles while recirculating traffic drains in 782 cycles with zero deadlocks observed. Under synthetic liquid traffic, transaction latency remains tightly bounded: mean latency is 4.87~$\mu$s (486.9 cycles), median is 3.86~$\mu$s (386 cycles), and P95 latency is 7.04~$\mu$s (704 cycles).

\subsection{Wordlength Precision Exploration}
\label{sec:eval_precision}
Evaluating fixed-point formats against the 740-DSP capacity of the Artix-7 200T reveals that narrower formats suffer severe numerical degradation: 24-bit Q6.18 exhibits an unacceptable MAE of 2,687~vol-bps due to dynamic range truncation when spot $S > \$64$, while Q7.21 incurs 1,067~vol-bps MAE from Horner polynomial underflow. Conversely, wider formats overrun device capacity: Q9.27 requires 848 DSPs (114.6\%), Q12.36 requires 1,152 DSPs (155.7\%), and FP32 requires 960 DSPs (129.7\%). Q8.24 represents the optimal operating point (608 DSPs, 82.2\%) satisfying both physical device limits and market accuracy requirements.

\begin{figure}[t]
\centering
\includefigfile[width=\columnwidth]{fig5_boundary_stress_surface.pdf}{figures/fig5_boundary_stress_surface.pdf}
\caption{Boundary-focused numerical stress surface and historical SPX quote replay: (a) 2,500-point 2D parameter error surface across moneyness $S/K \in [0.70, 1.40]$ and maturity $T \in [0.0027, 2.0]$~yr, identifying low-Vega ($T < 14$~days) and Pad\'{e} transition boundaries ($|S/K-1| \approx 0.15$); (b) historical replay of 1,382 recorded CBOE SPX quotes validating dynamic-range normalization with zero arithmetic overflow observed.}
\label{fig:boundary_stress}
\end{figure}

\subsection{Boundary Stress Surface \& Historical SPX Replay}
Figure~\ref{fig:boundary_stress} examines a 2,500-point boundary surface alongside offline market-data replay.

\textit{Boundary Stress Surface (Dataset C):} Moneyness $S/K \in [0.70, 1.40]$ and maturity $T \in [0.0027, 2.0]$~yr are sampled on a dense $50 \times 50$ grid. The surface identifies two demanding regimes:
\begin{enumerate}
    \item \textbf{Low-Vega Boundary ($T < 14$~days):} As $T \to 0$, Vega vanishes ($\nu \to 0$), causing unconstrained Newton-Raphson updates to destabilize. Hardware step clamping ($|\Delta\sigma| \le 0.25$) successfully bounds solver steps, maintaining finite error (MAE 4,389.72~vol-bps) without pipeline divergence.
    \item \textbf{Pad\'{e} Transition Zone ($|S/K - 1| \approx 0.15$):} Switching between the [1/1] Pad\'{e} approximant and CORDIC logarithm introduces a localized transition step (MAE 415.24~vol-bps, averaging 4.96 passes), reflecting an architectural trade-off of piecewise logarithm evaluation. Inside the grid interior ($0.85 \le S/K \le 1.15, T \ge 14\text{d}$), grid MAE is 104.38~vol-bps.
\end{enumerate}

\textit{Empirical SPX Market-Data Replay (Dataset D):} Replay of 1,382 recorded CBOE European SPX Weekly call quotes (recorded September 21, 2026; spot up to $\$7{,}643.99$, strike up to $\$7{,}580$, 15 expirations) serves as an out-of-sample stress test for dynamic-range normalization. Market prices were computed as midpoint quotes $C_{\text{mkt}} = (\text{bid} + \text{ask}) / 2$, time-to-maturity $T$ computed from calendar days, and risk-free benchmark $r = 5.25\%$. Scale normalization ($S^* = S/K, K^* = 1.0, C^* = C/K$) mapped four-digit spot prices into compact Q8.24 ($S^* \le 1.3711$) with zero arithmetic overflow observed across all 1,382 quotes. For NTM quotes, 72.1\% achieved single-pass convergence, with 437.93~vol-bps MAE (78.88~vol-bps median; liquid-moneyness MAE 487.62~vol-bps, 122.84~vol-bps median). The deep-wing subset ($N=15$, MAE 331.52~vol-bps) exhibited bounded solver behavior under hardware step clamping.

\begin{table}[t]
\caption{End-to-End RTL Validation Matrix Across Diverse Datasets}
\label{tab:rtl_validation}
\centering
\resizebox{\columnwidth}{!}{%
\begin{tabular}{lrrrr}
\toprule
\textbf{Validation Metric} & \textbf{Dataset A} & \textbf{Dataset B} & \textbf{Dataset C} & \textbf{Dataset D} \\
& \textbf{(Canonical 10k)} & \textbf{(Random 100k)} & \textbf{(Boundary 2.5k)} & \textbf{(SPX 1.38k)} \\
\midrule
\multicolumn{5}{l}{\textit{Transaction Integrity \& Execution Bounds}} \\
Evaluated Contracts ($N_{\text{acc}}$) & 10,000 & 100,000 & 2,500 & 1,382 \\
Converged Contracts ($N_{\text{conv}}$) & 9,673 (96.7\%) & 97,083 (97.1\%) & 1,127 (45.1\%) & 1,382 (100.0\%) \\
Convergence Failures ($N_{\text{fail}}$) & 327 (3.3\%) & 2,917 (2.9\%) & 1,373 (54.9\%) & 0 (0.0\%) \\
Retired Contracts ($N_{\text{ret}}$)   & 10,000 & 100,000 & 2,500 & 1,382 \\
Lost Contracts ($N_{\text{lost}}$)       & \textbf{0} & \textbf{0} & \textbf{0} & \textbf{0} \\
Duplicate Outputs ($N_{\text{dup}}$)    & \textbf{0} & \textbf{0} & \textbf{0} & \textbf{0} \\
TID Routing Mismatches                 & \textbf{0} & \textbf{0} & \textbf{0} & \textbf{0} \\
Arithmetic Overflows                   & \textbf{0} & \textbf{0} & \textbf{0} & \textbf{0} \\
\midrule
\multicolumn{5}{l}{\textit{Implied Volatility Error (Synthesized RTL vs. Golden Software)}} \\
Liquid Moneyness MAE (vol-bps)         & 7.67 & 14.84 & 2,702.22 & 487.62 \\
Overall Suite MAE (vol-bps)            & 127.86 & 133.87 & 6,996.26 & 485.92 \\
Overall Suite Median (vol-bps)         & 1.43 & 1.81 & 500.00 & 126.27 \\
Liquid Subset Median (vol-bps)         & 0.91 & 1.17 & 246.91 & 122.84 \\
Root-Mean-Square Error (RMSE, vol-bps) & 437.88 & 509.09 & 13,731.72 & 862.95 \\
95th Percentile Error (P95, vol-bps)   & 867.70 & 854.63 & 35,000.00 & 2,126.69 \\
99th Percentile Error (P99, vol-bps)   & 2,214.71 & 2,537.68 & 47,500.00 & 2,220.58 \\
Max Absolute Error ($|\Delta\sigma|$)  & 0.5549 & 0.8784 & 4.7500\textsuperscript{*} & 0.2266 \\
\midrule
\multicolumn{5}{l}{\textit{Concurrent Greeks Accuracy (MAE vs. Golden Software)}} \\
Delta MAE ($|\Delta_{\text{RTL}} - \Delta_{\text{gold}}|$) & 0.0119 & 0.0112 & 0.1651 & 0.0736 \\
Vega MAE ($|\nu_{\text{RTL}} - \nu_{\text{gold}}|$)       & 0.6296 & 0.6359 & 8.8059 & 0.0139\textsuperscript{\dag} \\
Gamma MAE ($|\Gamma_{\text{RTL}} - \Gamma_{\text{gold}}|$) & 0.0031 & 0.0032 & 0.0660 & 13.4204\textsuperscript{\ddag} \\
\bottomrule
\end{tabular}%
}
\vskip 1pt\raggedright
\footnotesize{\textsuperscript{*}Dataset C boundary errors reflect hardware clamping ($|\Delta\sigma| \le 0.25$ per step); across 8 passes, maximum cumulative error reaches 4.75 without numerical divergence. \\
\textsuperscript{\dag}Normalized Vega MAE ($\nu^*$). Standard monetary Vega MAE is 104.81 ($\nu = K\nu^*$). \\
\textsuperscript{\ddag}Dataset D Gamma MAE on the packed 26-bit egress bus reflects truncation of normalized $\Gamma^*$ for short expiries ($T < 14$d), where $\Gamma^*$ reaches 72.43 ($>1.999$ on 1,224 quotes); standard monetary Gamma $\Gamma = \Gamma^* / K$ is strictly $<0.010$ (mean $0.00175$). Across all 113,882 evaluated transactions, $N_{\text{acc}} = N_{\text{ret}}$ with 0 lost or duplicate contracts.}
\end{table}

\subsection{End-to-End Multi-Dataset RTL Validation}
Table~\ref{tab:rtl_validation} reports comprehensive validation across four datasets totaling 113,882 transactions: Canonical ($N=10\text{k}$), Randomized Stability ($N=100\text{k}$), Boundary Stress ($N=2.5\text{k}$), and SPX quotes ($N=1{,}382$). All 113,882 contracts retired without lost outputs ($N_{\text{lost}}=0$), duplicate emissions ($N_{\text{dup}}=0$), or TID mismatches. 

Numerical convergence is reported separately from transaction integrity: 96.7\% of Dataset~A, 97.1\% of Dataset~B, and 100.0\% of Dataset~D converged within the 8-pass budget. In Dataset~C, where deep out-of-the-money options have theoretical zero price, 45.1\% converged while 54.9\% retired gracefully at the 8-pass ceiling under hardware step clamping.

\begin{table}[t]
\caption{Full-System Saturation, Latency, and Robustness Diagnostics}
\label{tab:saturation_diagnostics}
\centering
\resizebox{\columnwidth}{!}{%
\begin{tabular}{lrrrrrr}
\toprule
\textbf{Workload /} & \textbf{Peak Ingress} & \textbf{Closed-Loop RTL} & \multicolumn{3}{c}{\textbf{Latency ($\mu$s / Cycles)}} & \textbf{Max In-Flight} \\
\textbf{Condition}  & \textbf{Capacity (MOps/s)} & \textbf{Rate (MOps/s)\textsuperscript{\dag}} & \textbf{Min} & \textbf{Median} & \textbf{P95} & \textbf{Ctx / Core} \\
\midrule
Dataset A (10k Canonical)  & 400.00 & 8.14 & 2.30 / 230 & 5.53 / 553 & 15.18 / 1,518 & 16 / 60 \\
Dataset B (100k Stability) & 400.00 & 8.55 & 2.30 / 230 & 5.52 / 552 & 15.18 / 1,518 & 16 / 60 \\
Dataset C (Boundary Grid)  & 400.00 & 5.34 & 2.30 / 230 & 15.18 / 1,518 & 15.21 / 1,521 & 16 / 60 \\
Dataset D (SPX Quotes)     & 400.00 & 18.42 & 2.30 / 230 & 2.30 / 230 & 5.55 / 555 & 16 / 60 \\
\midrule
\multicolumn{7}{l}{\textit{Downstream Backpressure Robustness (Dataset A, Randomized $m\_axis\_tready$ Stalls)}} \\
0\% Backpressure (Baseline) & 400.00 & 8.14 & 2.30 / 230 & 5.53 / 553 & 15.18 / 1,518 & 16 / 60 \\
20\% Backpressure           & 400.00 & 8.14 & 2.30 / 230 & 5.54 / 554 & 15.18 / 1,518 & 16 / 60 \\
50\% Backpressure           & 400.00 & 8.13 & 2.30 / 230 & 5.57 / 557 & 15.19 / 1,519 & 16 / 60 \\
\midrule
\multicolumn{7}{l}{\textit{Exp 1B Saturated Memory Ablation ($N=10{,}000$ contracts at 100 MHz signoff)}} \\
Zero-BRAM Engine (Proposed) & 400.00 & 8.18 & 2.30 / 230 & 5.53 / 553 & 15.18 / 1,518 & 0 BRAM \\
BRAM Baseline               & 400.00 & 8.14 & 2.31 / 231 & 5.56 / 556 & 15.27 / 1,527 & 28 BRAM18 \\
\bottomrule
\end{tabular}%
}
\vskip 1pt\raggedright
\footnotesize{\textsuperscript{\dag}Closed-loop RTL testbench throughput under a 16-credit/core window ($N=64$ credits array-wide; Little's Law bound $T = 64 / 6.984\,\mu\text{s} \approx 9.16$~MOps/s). Peak open-loop nominal datapath service capacity is 151.82~MOps/s for liquid traffic ($400 / 2.635$ passes). All configurations satisfy the $\le 60$ context headroom invariant.}
\end{table}

\subsection{Saturation Diagnostics \& Backpressure Robustness}
Table~\ref{tab:saturation_diagnostics} reports system saturation diagnostics. In closed-loop testbench co-simulation under a 64-credit window, the engine sustains 8.14~MOps/s on Dataset~A and 8.55~MOps/s on Dataset~B. Under downstream backpressure sweeps of 0\%, 20\%, and 50\% randomized stalls, the 64-entry distributed output FIFO absorbs stalls with $<0.2\%$ throughput variation (8.14~MOps/s at 0\% and 20\%, 8.13~MOps/s at 50\%) and only a 1.35-cycle mean latency delta (698.37 to 699.72 cycles).

\begin{table}[t]
\caption{Architectural Comparison with Prior FPGA Option Accelerators}
\label{tab:prior_art}
\centering
\resizebox{\columnwidth}{!}{%
\begin{tabular}{llcccccl}
\toprule
\textbf{Design} & \textbf{Numerical Task} & \textbf{Var. Iter.} & \textbf{Persist. Ctx.} & \textbf{Zero-BRAM} & \textbf{Greeks} & \textbf{Reported Rate / Clock} & \textbf{Target Device} \\
\midrule
Wang et al.~\cite{wang2022fpga} (2022) & Multi-Section IV & No & No & No (32 BRAM) & No & 120 MOps/s @ 120 MHz & Kintex UltraScale \\
Becker/Jin et al.~\cite{becker2011dynamic,jin2012optimising} (2011--12) & Finite Difference & No & Not rep. & No (Consumed) & No & 85--110 MOps/s & Virtex-6 LX240T \\
Tse et al.~\cite{tse2010customising} (2010) & Forward Pricing & No & Not rep. & No (Consumed) & No & 150 MOps/s @ 150 MHz & Virtex-4 FX100 \\
Klaisoongnoen et al.~\cite{klaisoongnoen2024evaluating} (2024) & Forward MC / AIE & No & Not rep. & No (Tile Mem) & No & Vector SIMD Engine & Versal VCK5000 \\
\midrule
\textbf{This Work (Artix-7)} & \textbf{Iterative IV + Greeks} & \textbf{Yes} & \textbf{Yes} & \textbf{Yes (0 BRAM)} & \textbf{Yes} & \textbf{400 MOps/s\textsuperscript{\dag} @ 100 MHz} & \textbf{Artix-7 200T} \\
\textbf{This Work (Alveo U50)} & \textbf{Iterative IV + Greeks} & \textbf{Yes} & \textbf{Yes} & \textbf{Yes (0 BRAM)} & \textbf{Yes} & \textbf{1.00 GOps/s @ 250 MHz} & \textbf{Alveo U50} \\
\bottomrule
\end{tabular}%
}
\vskip 1pt\raggedright
\footnotesize{\textsuperscript{\dag}Nominal one-pass contract-ingress capacity (151.82~MOps/s open-loop nominal liquid service capacity; 8.14~MOps/s closed-loop RTL testbench). AXI cold latency: 2.30~$\mu$s (Artix-7) / 0.920~$\mu$s (U50, 230 cyc); core: 2.27~$\mu$s / 0.908~$\mu$s (227 cyc). Prior works accelerate forward pricing~\cite{tse2010customising,klaisoongnoen2024evaluating}, finite difference~\cite{becker2011dynamic,jin2012optimising}, or root-finding without Greeks~\cite{wang2022fpga}. This work combines variable-iteration recirculation, persistent distributed context, zero-BRAM core storage, and concurrent Greeks.}
\end{table}

\subsection{Architectural Taxonomy \& Prior Art Comparison}
Table~\ref{tab:prior_art} compares the architectural features of the proposed engine against prior FPGA accelerators. Prior works either accelerate forward pricing~\cite{tse2010customising}, evaluate finite-difference grids~\cite{becker2011dynamic,jin2012optimising}, or perform multi-section interval search without variable-iteration loopback or Greeks~\cite{wang2022fpga}. While Klaisoongnoen et al.~\cite{klaisoongnoen2024evaluating} evaluate AMD Versal AI Engines (AIE) for nested Monte Carlo forward simulations, our work addresses the orthogonal challenge of streaming inverse implied-volatility solving. Within the surveyed accelerators, this work uniquely combines variable-iteration recirculation, persistent distributed context, zero-BRAM core storage, and concurrent Greek extraction.

As contextual background, a 16-thread C++ host CPU baseline (Intel i5-12500H, AVX2, 4.50~GHz Turbo, 45~W TDP) evaluating identical iterative IV convergence ($|\Delta C| \le \$0.01$, max 8 passes) and Greek extraction across 2M contracts achieves 14.22~MOps/s. While GPUs achieve high aggregate throughput on large-batch forward evaluations (e.g., CUDA SDK Black-Scholes sample~\cite{cuda2023samples}), variable-iteration root-finding incurs warp divergence and PCIe transfer latencies ($\gg 1\,\mu$s). The primary architectural evidence of this paper remains within the FPGA domain: the zero-BRAM context organization eliminates synchronous BRAM fetch read bubbles to achieve a +0.109~ns WNS advantage, enabling physical timing closure at 100~MHz.

% -------------------------------------------------------------------------
\section{Conclusion}
% -------------------------------------------------------------------------
We presented a four-core zero-BRAM spatial architecture for variable-iteration Black-Scholes implied-volatility solving and concurrent Greek extraction ($\Delta, \nu, \Gamma$). By combining Brenner-Subrahmanyam analytical seeding, Q8.24 fixed-point arithmetic, and distributed LUTRAM/SRL context storage, the engine achieves nominal post-route timing closure at 100.00~MHz on Artix-7 200T (400.00~MOps/s peak ingress, 4.258~W) and 250.00~MHz on Alveo U50 (1.00~GOps/s peak ingress, 9.560~W) with zero Block RAM or UltraRAM in the core. Priority loopback with a 4-slot scoreboard headroom invariant ($N_{\text{active}} \le 60$) prevents circular buffer deadlock while converging in an empirical average of 2.635 passes for liquid options (151.82~MOps/s nominal liquid service capacity on Artix-7; 8.14~MOps/s closed-loop RTL testbench throughput under a 64-credit window). End-to-end co-simulation across 113,882 transactions verifies zero transaction loss, 7.67~vol-bps liquid-moneyness IV MAE (median 0.91~vol-bps), and exact bit-level reproducibility against double-precision software models. Complete SystemVerilog RTL, Vivado implementation scripts, and DPI-C verification suites will be made openly available in an anonymized repository upon publication.

\begin{thebibliography}{10}

\bibitem{omahony2024role}
Aidan O'Mahony, Bernard Hanzon, and Emanuel Popovici. 2024.
\newblock The Role of {FPGAs} in Modern Option Pricing Techniques: A Survey.
\newblock \emph{Electronics} 13, 16 (2024), 3186.

\bibitem{black1973pricing}
Fischer Black and Myron Scholes. 1973.
\newblock The Pricing of Options and Corporate Liabilities.
\newblock \emph{Journal of Political Economy} 81, 3 (1973), 637--654.

\bibitem{brenner1988simple}
Menachem Brenner and Marti~G. Subrahmanyam. 1988.
\newblock A Simple Formula to Compute the Implied Standard Deviation.
\newblock \emph{Financial Analysts Journal} 44, 5 (1988), 80--83.

\bibitem{abramowitz1964handbook}
Milton Abramowitz and Irene~A. Stegun. 1964.
\newblock \emph{Handbook of Mathematical Functions with Formulas, Graphs, and Mathematical Tables}.
\newblock National Bureau of Standards.

\bibitem{cuda2023samples}
{NVIDIA Corporation}. 2023.
\newblock {CUDA} Samples: {Black-Scholes} Option Pricing.
\newblock \url{https://github.com/NVIDIA/cuda-samples}.

\bibitem{wang2022fpga}
Su~Wang, Hongxin Huan, Seng~Fat Wong, and Joseph Yen. 2022.
\newblock {FPGA} based Implied Volatility Calculation with Multi-section Method.
\newblock In \emph{Proc. IEEE 20th International Conference on Industrial Informatics (INDIN)}. 507--514.

\bibitem{becker2011dynamic}
Tobias Becker, Qiwei Jin, Wayne Luk, and Stephen Weston. 2011.
\newblock Dynamic Constant Reconfiguration for Explicit Finite Difference Option Pricing.
\newblock In \emph{Proc. International Conference on Reconfigurable Computing and FPGAs (ReConFig)}. 176--181.

\bibitem{jin2012optimising}
Qiwei Jin, Tobias Becker, Wayne Luk, and David~B. Thomas. 2012.
\newblock Optimising explicit finite difference option pricing for dynamic constant reconfiguration.
\newblock In \emph{Proc. 22nd International Conference on Field Programmable Logic and Applications (FPL)}. 165--172.

\bibitem{tse2010customising}
Arthur H.~T. Tse, David~B. Thomas, and Wayne Luk. 2010.
\newblock Customising Option Pricing: A Reconfigurable Core.
\newblock In \emph{Proc. 6th International Symposium on Applied Reconfigurable Computing (ARC)}. 24--35.

\bibitem{klaisoongnoen2024evaluating}
M.~Klaisoongnoen, N.~Brown, T.~Dykes, J.~R. Jones, and U.-U. Haus. 2024.
\newblock Evaluating {Versal AI} Engines for Option Price Discovery in Market Risk Analysis.
\newblock In \emph{Proc. 32nd ACM/SIGDA International Symposium on Field-Programmable Gate Arrays (ISFPGA)}. 176--182.

\end{thebibliography}

\end{document}
